\documentclass[10pt,a4paper]{article}
\usepackage{amsmath,amssymb,amsthm,graphicx}
\newtheorem{theorem}{Theorem}
\newtheorem{lemma}[theorem]{Lemma}
\newtheorem{definition}[theorem]{Definition}
\newtheorem{proposition}[theorem]{Proposition}
\newtheorem{corollary}[theorem]{Corollary}
\newtheorem{claim}[theorem]{Claim}
\newcommand{\XorEq}{\operatorname{XorEq}}
\setcounter{MaxMatrixCols}{15}
\title{NP-hardness of approximating constant-round total Boolean protocol leaf count}
\author{Shuichi Hirahara, Rahul Ilango, Bruno Loff}
\date{doi:10.4230/LIPIcs.CCC.2021.31; CC BY4.0}
\begin{document}
\maketitle
\noindent\textit{Editorial note.} Original authors: Shuichi Hirahara, Rahul Ilango and Bruno Loff. CCC2021 article31, 2021-07-08; DOI10.4230/LIPIcs.CCC.2021.31.\par\medskip
\noindent\textit{Editorial note.} License: https://creativecommons.org/licenses/by/4.0/\par\medskip
\noindent\textit{Editorial note.} Faithful editable transcription from the official CCC2021 CC BY4.0 published version, article31, 2021-07-08. Theorem4 and its formal Corollary24 conclusion are one constant-round deterministic protocol-size hardness result. Full three-round Boolean-matrix construction, chromatic-number gap prerequisite, large-alphabet warmup, cost characterization, deterministic fooling-set round elimination and approximation reduction proofs are retained with original numbered statements. Original total Boolean NxN input, constant round/approximation conditions, Alice/Bob-first conventions and model definitions are retained. Original mathematical expressions and reduction-time wording are transcribed without author-proof correction or added assumptions. The complementary leaf-set overlines were read from official p31:18; typography-only line breaks preserve all terms. All65 official bibliography entries retained. Randomized/conditional extensions in Sections6--7 are not claimed by this deterministic item; the full official original remains available. Supporting statements and direct consequences are not counted separately.\par\medskip
\noindent\textit{Allowed standard prerequisites (editorial).}\begin{itemize}
\item Chromatic-number hardness of approximation. Exact Theorem9 with original quantifiers and references[42,24,20,65] retained.
\item Explicit universal sets. Exact Definition15 and Lemma16[51], including size/time parameters, retained; no recursion into prior proof.
\end{itemize}\medskip
\setcounter{theorem}{3}
\begin{theorem}
For any $d\ge3$ there exists a constant $\Delta_d>0$ such that the following holds. If there exists a polynomial-time algorithm which, when given a total two-player Boolean-valued function $f:[N]\times[N]\to\{0,1\}$ represented as a Boolean matrix of dimensions $N\times N$, approximates $L_d^A(f)$ within a factor of $1+\Delta_d$, then $\mathrm P=\mathrm{NP}$.
\end{theorem}

\section*{Author preliminaries and communication model (Section 2)}
For a positive integer $n$, we let $[n]$ denote the set $\{1,\ldots,n\}$.
\paragraph{Binomial Coefficients and Projections.} The binomial coefficient $\binom nk$ equals the number of distinct subsets of $[n]$ with exactly $k$ elements. Similarly, $\binom n{\le k}$ denotes the number of distinct subsets of $[n]$ that have at most $k$ elements. Finally, $\binom{[n]}k$ denotes the set of all subsets of $[n]$ with exactly $k$ elements.
If $x\in\{0,1\}^n$ and $S=\{i_1,\ldots,i_k\}\in\binom{[n]}k$, then $x_S\in\{0,1\}^k$ denotes the projection of $x$ to $S$, given by $x_S=x_{i_1}\cdots x_{i_k}$.
\paragraph{Communication complexity.} We assume basic familiarity with communication complexity \cite{ref38}. We write the definitions here for clarity.
\setcounter{theorem}{4}
\begin{definition}[Protocol]\label{def:protocol}
Let $\mathcal A,\mathcal B,\mathcal Z$ be finite sets. A deterministic protocol $\pi$ over $\mathcal A\times\mathcal B\times\mathcal Z$ is a rooted tree:
\begin{itemize}
\item Each node $v$ is associated with a rectangle $\pi^{-1}(v)=A\times B$, with $A\subseteq X$ and $B\subseteq Y$.
\item Each non-leaf node $v$, associated with a rectangle $\pi^{-1}(v)=A\times B$, is labeled by either (a) a partition $A=\mathop{\dot\bigcup}_{c\in P_v} A_c$ of $A$, in which case we say it is Alice's node or (b) a partition $B=\mathop{\dot\bigcup}_{c\in P_v} B_c$ of $B$, in which case we say it is Bob's node.
\item Each leaf node is labeled by an element of the output domain $\mathcal Z$.
\item The rectangle associated with the root is the input domain $\mathcal A\times\mathcal B$.
\item If a non-leaf node $v$ of Alice is associated with rectangle $\pi^{-1}(v)=A\times B$ and (a) labeled by a partition $A=\mathop{\dot\bigcup}_{c\in P_v} A_c$ of $A$, then for each $c\in P_v$ there will be one child $v_c$ of $v$, which will be associated with the rectangle $\pi^{-1}(v_c)=A_c\times B$; similarly for Bob's nodes.
\end{itemize}
We let the leaf complexity of $\pi$, written $L(\pi)$, be the number of leaf nodes of $\pi$. We let the round complexity of $\pi$, written $R(\pi)$, be the height of $\pi$, i.e., the maximum number of edges in any root-to-leaf path of $\pi$.
A root-to-leaf path $v_1\to\cdots\to v_{k+1}$ is said to have communication length $\sum_{i=1}^k\lceil\log|P_{v_i}|\rceil$ (where $|P_{v_i}|$, as defined above, is the number of parts in the partition of $A$ or $B$ associated with node $v_i$). We then let the communication complexity of $\pi$, written $C(\pi)$, be the maximum communication length of any root-to-leaf path of $\pi$.
Given $(a,b)\in\mathcal A\times\mathcal B$, we let $\pi(a,b)$ denote the (unique) leaf $v$ of $\pi$ having $(a,b)\in\pi^{-1}(v)$. For $z\in\mathcal Z$, we may write $\pi(a,b)=z$ to mean that the leaf $\pi(a,b)$ is labeled by $z$.
A randomized protocol over $\mathcal A\times\mathcal B\times\mathcal Z$ is a distribution over deterministic protocols over $\mathcal A\times\mathcal B\times\mathcal Z$. We will use a boldface Greek letter, such as $\boldsymbol\pi$, to denote a protocol sampled from this distribution. We then let $L(\boldsymbol\pi)$ be the maximum $L(\pi)$ over all $\pi$ in the support of $\boldsymbol\pi$, and likewise for $R(\boldsymbol\pi)$ and $C(\boldsymbol\pi)$.
\end{definition}
\begin{definition}
A function $f:\mathcal A\times\mathcal B\to\mathcal Z$ is said to be computed by a deterministic protocol $\pi$ over $\mathcal A\times\mathcal B\times\mathcal Z$ if we have $f(a,b)=\pi(a,b)$ for every $(a,b)\in\mathcal A\times\mathcal B$. Furthermore, if $\varepsilon\in[0,1]$, then $f$ is said to be computed with error $\varepsilon$ by a randomized protocol $\boldsymbol\pi$ if, for every $(a,b)\in\mathcal A\times\mathcal B$, $\Pr[\boldsymbol\pi(a,b)=f(a,b)]>1-\varepsilon$ (the probability is over the choice of $\boldsymbol\pi$).
We may then define:
\begin{itemize}
\item $L_d^A(f)$ is the minimum leaf complexity $L(\pi)$ among all deterministic protocols $\pi$ that compute $f$, and have round complexity $R(\pi)\le d$, and such that the root node of $\pi$ is Alice's. $L_d^B(f)$ is defined likewise, but for protocols where the root node is Bob's.
\item $L_{d,\varepsilon}^A(f)$ is the minimum $L(\boldsymbol\pi)$ among all randomized protocols $\boldsymbol\pi$ that compute $f$ with error $\varepsilon$, and have $R(\boldsymbol\pi)\le d$, and such that the root node of $\boldsymbol\pi$ is (always) Alice's. $L_{d,\varepsilon}^B(f)$ is defined likewise for randomized protocols where the root node is (always) Bob's.
\item $C_d^A,C_d^B,C_{d,\varepsilon}^A,C_{d,\varepsilon}^B$ are defined analogously where the communication complexity $C(\pi)$ replaces the leaf complexity $L(\pi)$.
\end{itemize}
\end{definition}
Since the communication transcript of a given run of protocol determines the leaf, it must follow that:
\begin{proposition} For any protocol $\pi$, $\log L(\pi)\le C(\pi)$.\end{proposition}
\paragraph{Chromatic number.} All our NP-hardness reductions are from the chromatic number problem:
\begin{definition}[Chromatic number] A coloring of an undirected graph $G$, is a partition of the vertices such that no edge has both endpoints in the same part. The chromatic number of a graph $G$, denoted $\chi(G)$, is the smallest number of parts in a coloring of $G$.\end{definition}
The NP-hardness of approximating the chromatic number has been established by a series of results \cite{ref42,ref24,ref20}, culminating in a paper by Zuckerman \cite{ref65}, where the following was proven:
\begin{theorem}[Hardness For Chromatic Number]\label{thm:chromatic}
For every $\epsilon>0$, there is a deterministic polynomial time algorithm that on an input $x\in\{0,1\}^*$ outputs a graph $G$ on $n$ vertices such that
\begin{itemize}
\item if $x$ is a YES instance of SAT, then $\chi(G)\le n^\epsilon$, and
\item if $x$ is a NO instance of SAT, then $\chi(G)\ge n^{1-\epsilon}$.
\end{itemize}
\end{theorem}
\section*{Author three-round cost characterization (Section 3)}
We start by observing that deterministic 3-round communication complexity may be approximated by a very simple combinatorial quantity.
\begin{definition}
Let $A,B,Z$ be finite sets, and let $M$ be an $A\times B$ matrix over $Z$. Let $\mathcal P=\{P_i\}_{i\in[k]}$ be a partition of (the columns) $B$ for some $k\in\mathbb N$. (That is, $\emptyset\ne P_i\subseteq B$, $\bigcup_{i\in[k]}P_i=B$ and $P_i\cap P_j=\emptyset$ for every $i\ne j$.) For a subset $P\subseteq B$ of columns, we denote by $\operatorname{Cost}_M(P)$ the number of distinct rows of $M$ restricted to columns in $P$, i.e.,
\[\operatorname{Cost}_M(P)=|\{x_P\in Z^P\mid x\in Z^B\text{ is a row of }M\}|.\]
We further define $\operatorname{Cost}_M(\mathcal P)$ to be $\sum_{i=1}^k\operatorname{Cost}_M(P_i)$.
\end{definition}
\begin{proposition}\label{prop:cost}
Let $f:A\times B\to Z$ be a function and $M$ be the $A\times B$ communication matrix (with entries in $Z$) that corresponds to $f$. Let $q\in\mathbb N$ be the maximum number of distinct values in a single row of $M$. We then have
\[ L\le L_3^B(f)\le L\cdot q,\qquad\text{where }L=\min_{\mathcal P}\operatorname{Cost}_M(\mathcal P),\]
and where the minimum is taken over all partitions $\mathcal P$ of $B$. Furthermore we have that
\[\log L\le C_3^B(f)\le\log L+\log q+O(1).\]
\end{proposition}
\begin{proof}
It may be easily seen that $L_2^A(f)$ is lower-bounded by the number of distinct rows in the communication matrix of $f$. Because if Alice's partition of the rows includes a part with two different rows, then Bob's ensuing partition of the columns cannot avoid having a non-monochromatic column. It then follows that, if we have a 3-round protocol $\pi$ where Bob speaks first, and $\mathcal P$ is Bob's partition of the columns in the first round, then $\operatorname{Cost}_M(\mathcal P)$ is a lower-bound on smallest number of leaves that $\pi$ needs to use to compute $f$, and thus $L_3^B(f)\ge L$. The lower-bound on $C_3^B(f)$ now follows from Proposition~7.
Conversely, it is also easy to see that $L_2^A(f)$ is upper-bounded by $q$ times the number of distinct rows in the communication matrix of $f$. The protocol that achieves this bound has Alice tell which kind of row she has, and now in each rectangle the rows are all equal, hence Bob can just tell Alice the color of his column in this row, of which there are $q$ possibilities. Thus $L_3^B(f)\le L\cdot q$, and the same protocol shows that $C_3^B(f)\le\log L+\log q+O(1)$.
\end{proof}

In general, the approximate factor $q$ of Proposition~11 can be as large as $|Z|$. However, in the following construction, we will construct a matrix where each row has at most $q=3$ distinct values; in this case, Proposition~11 guarantees that $\min_{\mathcal P}\operatorname{Cost}_M(\mathcal P)$ provides a 3-factor approximation of $L_3(M)$.
\begin{theorem}
Given an undirected graph $G=([n],E)$ with $n$ vertices and $|E|=m>0$ edges, one may construct in deterministic polynomial time a function $f_G:[a]\times[n]\to\{0,1,\ldots,\ell\}$, with $\ell=m^2n$, $k=\sqrt{n\ell}=mn$ and $a=\ell+m\cdot k^2$, such that
\[\chi(G)\cdot\ell\le L_3^B(f_G)\le\chi(G)\cdot6\ell.\]
Furthermore, we also have
\[\log\chi(G)+\log\ell\le C_3^B(f_G)\le\log\chi(G)+\log\ell+O(1).\]
\end{theorem}
\begin{proof}
We let $A_\ell$ denote the $\ell\times1$ column vector,
\[A_\ell=\begin{bmatrix}1\\\vdots\\\ell\end{bmatrix}.\]
We let $B_k$ and $C_k$ denote the $k^2\times1$ column vectors
\[B_k=\begin{bmatrix}1\\\vdots\\1\\\vdots\\k\\\vdots\\k\end{bmatrix},\qquad C_k=\begin{bmatrix}1\\\vdots\\k\\\vdots\\1\\\vdots\\k\end{bmatrix},\]
where each value $i\in[k]$ appears $k$ times. Given an edge $\{v,w\}\in[n]\times[n]$, with $v<w$, we define the $k^2\times n$ matrix $M$ so that
\[M_{\{v,w\}}=\begin{bmatrix}0&\cdots&0&B_k&0&\cdots&0&C_k&0&\cdots&0\end{bmatrix},\]
where $B_k$ appears in the $v$-th column and $C_k$ in the $w$-th column. Finally, let $E=\{e_1,\ldots,e_m\}$ denote the edges of $G$; then we define the $a\times n$ communication matrix $f_G$ so that
\[f_G=\begin{bmatrix}A_\ell&\cdots&A_\ell\\\multicolumn{3}{c}{M_{e_1}}\\\multicolumn{3}{c}{\vdots}\\\multicolumn{3}{c}{M_{e_m}}\end{bmatrix}\in\{0,1,\ldots,\ell\}^{a\times n},\]
where $a=\ell+m\cdot k^2$. Observe that each row of $f_G$ has at most $q=3$ distinct values; thus, Proposition~11 provides a 3-factor approximation. We now show the stated inequality, namely, that
\[\chi(G)\cdot\ell\le L_3^B(f_G)\le\chi(G)\cdot2q\ell.\]
The upper-bound is easy to see. Given a coloring of $G$ into $\chi(G)$ colors, we may take the 3-round protocol $\pi$ where Bob first tells Alice which color his vertex has. This partitions the columns by a partition $\mathcal P=\{P_c\}_{c\in[\chi(G)]}$ formed of the various color classes of our coloring. Fix any color $c$ and consider $\operatorname{Cost}_{f_G}(P_c)$. Since $P_c$ is an independent set of $G$, the $C_k$ and $B_k$ columns of each $M_{e_i}$ sub-matrix will always be placed in different parts; therefore, $\operatorname{Cost}_{f_G}(P_c)\le\ell+mk$, where the first term counts the number of rows in $A_\ell$ and the second term counts the number of distinct rows in $M_{e_i}$ restricted to columns of $P_c$ for each $i\in[m]$. Using Proposition~11, we obtain
\[L_3^B(f_G)/q\le\operatorname{Cost}_{f_G}(\mathcal P)=\sum_c\operatorname{Cost}_{f_G}(P_c)\le\chi(G)\cdot(\ell+mk)=\chi(G)\cdot2\ell,\]
and Proposition~11 also gives us
\[C_3^B(f_G)\le\log\chi(G)+\log\ell+O(1).\]
For the other direction, let $\pi$ be any 3-round protocol for $f_G$. The first round of $\pi$ partitions the columns of $f_G$, which is to say, it partitions the vertices of $G$ by some partition $\mathcal P=\{P_i\}_{i\in I}$. We claim that $\chi(G)\cdot\ell\le\operatorname{Cost}_{f_G}(\mathcal P)$ by analyzing the following two cases.
\begin{enumerate}
\item If the partition $\mathcal P$ does not form a coloring of $G$, then there must exist an edge $\{v,w\}$ such that the $v$-th column and the $w$-th column of $G$ are placed in the same part $P_i$. This means that the $C_k$ and $B_k$ columns of the $M_{\{v,w\}}$ sub-matrix are both placed in $P_i$. In this case, we have $\operatorname{Cost}_{f_G}(P_i)\ge k^2\ge n\cdot\ell\ge\chi(G)\cdot\ell$, and thus the lower bound holds.
\item Otherwise, suppose that the partition $\mathcal P$ does form a coloring of $G$. Then the number of parts is $\ge\chi(G)$, and so just the contribution from the first $\ell$ rows gives us $\operatorname{Cost}_{f_G}(\mathcal P)\ge\chi(G)\cdot\ell$.
\end{enumerate}
The lower-bounds on $L_3^B(f_G)$ and $C_3^B(f_G)$ then follow from Proposition~11.
\end{proof}
The following corollary shows that it is NP-hard to approximate $L_3^B(f)$, for a given total two player function $f:[N]\times[N]\to[N]$, with an approximation ratio better than (roughly) $N^{1/5}$.
\begin{corollary}
For every $L\subseteq\{0,1\}^*$ in $\mathrm{NP}$ and every constant $\varepsilon>0$, there exists a polynomial-time algorithm that, on input $x\in\{0,1\}^*$, outputs a communication matrix $M\in[N]^{N\times N}$ such that if $x\in L$ then $L_3^B(M)\le N$, and if $x\notin L$ then $L_3^B(M)>N^{6/5-\varepsilon}$.
\end{corollary}
\begin{proof}
We compose the hardness of approximation result for chromatic number in Theorem~9 with the reduction of Theorem~12, padding the communication matrix with all-0 columns to make it square (since Bob speaks first, this adds at most one leaf), so the communication matrix of $f$ is an $N\times N$ matrix with $N=a=\Theta(n^5)$. The parameters then come from the fact that $n=\Theta(N^{1/5})$.
\end{proof}
Using the bounds on $C_3^B$ instead of the bounds on $L_3^B$ from Theorem~12, we conclude that it is NP-hard to approximate $C_3^B(f)$, for a given total two player function $f:[N]\times[N]\to[N]$, with an approximation ratio better than (roughly) $6/5$. More precisely:
\begin{corollary}
For every $L\subseteq\{0,1\}^*$ in $\mathrm{NP}$ and every $\varepsilon>0$, there exists a polynomial-time algorithm that, on input $x\in\{0,1\}^*$, outputs a communication matrix $M\in[N]^{N\times N}$ such that if $x\in L$ then $C_3^B(M)\le\log N$, and if $x\notin L$ then $C_3^B(M)>(\frac65-\varepsilon)\log N$.
\end{corollary}

\section*{Hardness for deterministic 3-round protocols (author Section 4)}
Building on the proof ideas presented in Section 3, in this section, we prove NP-hardness of approximating the communication complexity of deterministic 3-round protocols. A key building block is to use the notion of universal set.
\setcounter{theorem}{14}
\begin{definition}[Universal set]
Let $r,c,k\in\mathbb N$, and let $M\in\{0,1\}^{r\times c}$ be a matrix whose columns are $M^{(1)},\ldots,M^{(c)}$. We say that $M$ is $(c,k)$-universal if, for every set $\{y_1,\ldots,y_k\}\subseteq[c]$ of $k$ columns of $M$, the matrix
\[\begin{bmatrix}M^{(y_1)}&\cdots&M^{(y_k)}\end{bmatrix}\]
has $2^k$ distinct rows. The set of all the rows of $M$ is called a $(c,k)$-universal set.
\end{definition}
We use the explicit construction of a universal set due to \cite{ref49}.
\begin{lemma}[Naor, Schulman and Srinivasan \cite{ref51}]
There exists a deterministic algorithm that, given $c$ and $k\in\mathbb N$, outputs a $(c,k)$-universal matrix $M\in\{0,1\}^{r\times c}$ such that $r=2^{k+O(\log^2 k)}\cdot\log c$ in time a polynomial in $c$ and $2^k$.
\end{lemma}
Now we state the main result of this section.
\begin{theorem}\label{thm:boolean-construction}
Let $\epsilon>0$ be an arbitrary constant. Given an undirected graph $G=([n],E)$ with $n$ vertices and $|E|=m>0$ edges, one may construct in deterministic polynomial time a function $f_G:[a]\times[b]\to\{0,1\}$ and a number $\ell\in\mathbb N$, with $a,b,\ell=O(n^8)$, such that
\[\chi(G)\cdot\ell\le L_3^B(f_G)\le\chi(G)\cdot\ell^{1+\epsilon}\]
and
\[\log\chi(G)+\log\ell\le C_3^B(f_G)\le\log\chi(G)+(1+\epsilon)\cdot\log\ell.\]
\end{theorem}
The idea of the proof is to build upon the reduction in the proof of Theorem~12, by replacing each column with entries from $[\ell]$ with a block of columns that have entries from $\{0,1\}$. The difficulty in making this work is that a partition of the columns might not respect our blocks and could place columns from the same block into different parts. We solve this by thinking as follows. Either the partition is large, meaning it has many parts, so the protocol also has many leaves, which proves our lower bound, or otherwise for any block $C_v$ of columns the part $P_i$ which has most columns from $C_v$ has many columns from $C_v$; we may then act as if ``$C_v$ was placed in $P_i$''.
\begin{proof}
Let $t,k,c\in\mathbb N$ be parameters chosen later. Let $A\in\{0,1\}^{r\times c}$ and $M\in\{0,1\}^{s\times c}$ be the $(c,t)$-universal and $(c,k)$-universal matrices, respectively, that are constructed by the polynomial-time deterministic algorithm of Lemma~16; then we have $r=2^{(1+o(1))\cdot t}\cdot\log c$ and $s=2^{(1+o(1))\cdot k}\cdot\log c$.
Let $x_1,\ldots,x_s$ be the rows of $M$. We then let $B$ and $C$ denote the $s^2\times c$ matrices
\[
B=\begin{bmatrix}x_1\\\vdots\\x_1\\\vdots\\x_s\\\vdots\\x_s\end{bmatrix},\qquad
C=\begin{bmatrix}x_1\\\vdots\\x_s\\\vdots\\x_1\\\vdots\\x_s\end{bmatrix},
\]
where a row $x_i$ appears $s$ times for each $i\in[s]$. Given an edge $\{v,w\}\in[n]\times[n]$, with $v<w$, we define the $s^2\times c\cdot n$ matrix:
\[M_{\{v,w\}}=\begin{bmatrix}0&\cdots&0&B&0&\cdots&0&C&0&\cdots&0\end{bmatrix},\]
where $B$ appears in the $v$-th block of $c$ columns and $C$ in the $w$-th block of $c$ columns. Finally, let $E=\{e_1,\ldots,e_m\}$ denote the edges of $G$; then we define the $(r+m\cdot s^2)\times c\cdot n$ communication matrix $f_G$ so that
\[f_G=\begin{bmatrix}A&\cdots&A\\\multicolumn{3}{c}{M_{e_1}}\\\multicolumn{3}{c}{\vdots}\\\multicolumn{3}{c}{M_{e_m}}\end{bmatrix}\in\{0,1\}^{(r+m\cdot s^2)\times c\cdot n}.\]
Let $\ell:=2^t$. Let $\mathcal P$ be a partition of the columns of $f_G$ that minimizes $\operatorname{Cost}_{f_G}(\mathcal P)$. We now show that, for a suitable choice of $t,k$ and $c$,
\[\chi(G)\cdot\ell\le\operatorname{Cost}_{f_G}(\mathcal P)\le\chi(G)\cdot\ell^{1+o(1)}.\]
This will complete the proof, because Proposition~11 shows that, for a binary matrix $f_G$, $\operatorname{Cost}_{f_G}(\mathcal P)$ is a 2-factor approximation of $L_3^B(f_G)$ and $\log\operatorname{Cost}_{f_G}(\mathcal P)$ is an approximation of $C_3^B(f_G)$ up to an additive $O(1)$ term. We will choose the parameters $t,k$ and $c$ so that the following conditions are satisfied.
\begin{enumerate}
\item Condition 1. $r+ms\le\ell^{1+o(1)}$.
\item Condition 2. $c\ge n\ell\cdot\max\{t,k\}$.
\item Condition 3. $2^{2k}\ge n\ell$.
\end{enumerate}
Given that Condition 1 is satisfied, the complexity upper-bounds are easy to see. Let $\mathcal P$ be a coloring of $G$ that partitions the vertex set into $\chi(G)$ parts. Since no class contains an edge, the $C$ and $B$ sub-matrices of each $M_{e_i}$ sub-matrix will always be placed in different parts, and we thus have
\[\operatorname{Cost}(\mathcal P)\le\chi(G)\cdot(r+ms)\le\chi(G)\cdot\ell^{1+o(1)}.\]
For the other direction, we will lower-bound $\operatorname{Cost}(\mathcal P)$ for every partition $\mathcal P$. To begin, if the number of parts is $|I|\ge n\ell$, then we must conclude that $\operatorname{Cost}(\mathcal P)\ge n\ell\ge\chi(G)\cdot\ell$, as desired.
Otherwise, $|I|\le n\ell$. Then for every block of columns $v\in[n]$ there must exist a part $P_{i(v)}$ which contains at least $c/|I|\ge c/n\ell$ columns from the $v$-th block. Observe that $c/n\ell\ge\max\{t,k\}$ by Condition 2.
Now, either the mapping $v\mapsto i(v)$ is a valid coloring of $G$ or not. First, suppose that the mapping $v\mapsto i(v)$ is not a valid coloring of $G$, meaning that there exists an edge $\{v,w\}\in E$ such that $i=i(v)=i(w)$. Then the $v$-th column block and the $w$-th column block each have at least $\max\{t,k\}$ columns in the same part $P_i$. This will mean that the $C$ and $B$ sub-matrices of the $M_{\{v,w\}}$ sub-matrix each have at least $k$ columns in $P_i$. But then, since $M$ is $(c,k)$-universal, it follows from Condition 3 that $\operatorname{Cost}(\mathcal P)\ge\operatorname{Cost}(P_i)\ge2^{2k}\ge n\ell$.
Next, suppose that $i:[n]\to\mathcal P$ does form a coloring of $G$. Then there exist at least $\chi(G)$ parts $P_i$ each receiving at least $t$ columns from some column block, and so, since $A$ is $(c,t)$-universal, the contribution from the $A$ columns give us $\operatorname{Cost}(\mathcal P)\ge\chi(G)\cdot2^t=\chi(G)\cdot\ell$.
This will give us the theorem, and all we are left to do is ensure that the various conditions can be met: We define $t$ so that $\ell=2^t$ satisfies that $\ell\le(nm^2)^{1+\epsilon/2}<2\ell$. To meet Condition 3, let $k$ be the smallest integer satisfying that $n\ell\le2^{2k}$. Let $c:=n\ell\cdot\max\{t,k\}$ so that Condition 2 is satisfied. Finally, observe that Condition 1 is satisfied because
\begin{align*}
r+ms&\le\ell^{1+o(1)}\cdot\log c+m\cdot2^{(1+o(1))\cdot k}\cdot\log c\\
&\le\ell^{1+o(1)}+m(n\ell)^{\frac12+o(1)}\cdot\log c\le\ell^{1+\epsilon},
\end{align*}
where the last inequality holds for all large $n,m$.
\end{proof}
We can now prove that $L_3^A$ and $C_3^A$ are hard to approximate, also for Boolean functions.
\begin{theorem}\label{thm:three-round-hardness}
For every constant $\varepsilon>0$, there exists a deterministic quasipolynomial-time algorithm that, on input $x\in\{0,1\}^*$, outputs a communication matrix $M\in\{0,1\}^{N\times N}$ and a number $k\in\mathbb N$, with $k\le N=|x|^{O(1)}$, such that
\begin{enumerate}
\item if $x$ is a YES instance of SAT, then $L_3^B(M)\le k$ and $C_3^B(M)\le\log k$, and
\item if $x$ is a NO instance of SAT, then $L_3^B(M)\ge N^{\frac18-\varepsilon}\cdot k$ and $C_3^B(M)\ge\log k+(\frac18-\varepsilon)\log N$.
\end{enumerate}
\end{theorem}
\begin{proof}
We combine the two reductions of Theorems~9 and~17, which we invoke with the same parameter $\varepsilon$. Fix any input $x$ and let $G$ be an $n$-vertex graph that is produced by the reduction of Theorem~9 on input $x$. Let $M\in\{0,1\}^{a\times b}$ be the communication matrix of $f_G$, and $\ell\in\mathbb N$, be given by the reduction of Theorem~17 on input $G$, and set $N=\max(a,b)=O(n^8)$, $k=n^\varepsilon\cdot\ell^{1+\varepsilon}$. We may assume that $a=b$ without loss of generality, since otherwise we may pad $M$ with all-0 rows or all-0 columns, and this changes the leaf complexity by at most a factor of 2, and the communication complexity by at most 1 bit, so this change makes no difference to the result.
We verify that the inequalities are satisfied for $M$: If $x\in L$, then we have $L_3^B(M)\le\chi(G)\cdot\ell^{1+\epsilon}\le n^\varepsilon\cdot\ell^{1+\epsilon}=k$. If $x\notin L$, then we have
\[L_3^B(M)\ge\chi(G)\cdot\ell\ge n^{1-\epsilon}\cdot\ell\ge N^{\frac18-O(\epsilon)}\cdot k.\]
Similar inequalities hold for $C_3^B(M)$. Since $\varepsilon$ can be arbitrarily small, we may ignore the constant factors.
\end{proof}

\section*{From 3-rounds to multiple rounds using deterministic round elimination (author Section 5)}
In this section we show how we can use an oracle for computing $L_{d+1}^B$ in order to compute $L_d^A$. The gadget in our reduction involves a special function, $\XorEq_k$, which is a small modification of the standard Equality function.
\begin{definition}
The function $\XorEq_k:[k]\times([k]\times\{0,1\})\to\{0,1\}$ is given by
\[\XorEq_k(x;y,z)=\begin{cases}z&\text{if }x\ne y,\\1-z&\text{if }x=y.\end{cases}\]
\end{definition}
The key property of $\XorEq_k$ is that it has a fooling set lower bound that is tight. In particular, $\{(x;y,z):x=y\}$ is a fooling set of cardinality $2k$, and there is a 2-round protocol for solving $\XorEq_k$ with $2k$ leaves (where Alice just sends her full input to Bob, and he replies with the output).
\begin{lemma}[Fooling set lower-bound for $\XorEq_k$]\label{lem:fooling}
Let $\pi$ be a protocol for solving the function $f:([k]\times[a])\times([k]\times\{0,1\})\to\{0,1\}$ given by $f(x_0,x_1;y,z)=\XorEq_k(x_0;y,z)$.\footnote{In our definition of $f$, the input $x_1$ does not affect the output of the function. The fact that this lemma holds even when there is an extraneous input like $x_1$ will be used later.} Then
\[\pi(x_0,x_1;y,z)\ne\pi(x'_0,x'_1;y',z')\]
if either
\begin{itemize}
\item $y=x_0\ne x'_0$, or
\item $y=x_0=x'_0$ and $z\ne z'$.
\end{itemize}
\end{lemma}
\begin{proof}
First, suppose for contradiction that $\pi(x_0,x_1;y,z)=\pi(x'_0,x'_1;y',z')$ and $y=x_0\ne x'_0$. Since leaves are combinatorial rectangles, we can infer that $\pi(x'_0,x'_1;y,z)=\pi(x_0,x_1;y,z)$. But since $y=x_0\ne x'_0$, we know that
\[f(x_0,x_1;y,z)=1-z\ne z=f(x'_0,x'_1;y,z)\]
so this contradicts that $\pi(x_0,x_1;y,z)$ is a monochromatic leaf.
Similarly, if $y=x_0=x'_0$ and $z\ne z'$, then we have $f(x_0,x_1;y,z)\ne f(x'_0,x'_1;y',z')$, so $\pi(x_0,x_1;y,z)\ne\pi(x'_0,x'_1;y',z')$ by monochromaticness.
\end{proof}
The main technical portion of our reduction is the following deterministic variant of the round-elimination lemma.
\begin{lemma}[Restatement of Lemma 3]\label{lem:round-elimination}
Let $d\ge3$. Let $f:[a]\times[b]\to\{0,1\}$. Let $F:([k]\times[a])\times([k]\times[b]\times\{0,1\}\times\{0,1\})\to\{0,1\}$ be given by
\[F(x_0,x_1;y_0,y_1,z,i)=\begin{cases}\XorEq_k(x_0;y_0,z)&\text{if }i=0,\\f(x_1;y_1)&\text{if }i=1.\end{cases}\]
Then we have
\[\min\{4k,2k-2+L_d^B(f)\}\le L_{d+1}^A(F)\le2k+L_d^B(f).\]
\end{lemma}
\begin{proof}
The upper bound comes from the protocol where Alice skips the first round of communication, Bob sends $i$ to Alice and begins running the best $d$-round Bob-first protocol for $\XorEq_k$ or $f$, based on whether $i=0$ or $i=1$. In particular, we have that
\[L_{d+1}^A(F)\le L_d^B(\XorEq_k)+L_d^B(f)\le2k+L_d^B(f),\]
where the last upper bound uses that $d\ge3$.
We now argue the lower bound. Suppose $\pi$ is a $(d+1)$-round Alice-first protocol for $F$. Let $\mathcal L=\{\ell_1,\ldots,\ell_{L(\pi)}\}$ denote the set of leaves of $\pi$.
Our arguments split into two cases depending on whether there is a good input $x_0\in[k]$. We say an input $x_0\in[k]$ is good if all of Alice's inputs that begin with $x_0$ are placed in a single partition. More formally, let $\mathcal P=\{P_1,\ldots,P_{|\mathcal P|}\}$ be the partition of Alice's inputs corresponding to the first round of $\pi$. We say $x_0\in[k]$ is good if there exists a $q\in[|\mathcal P|]$ such that $\{x_0\}\times[a]\subseteq P_q$.
\paragraph{Case 1: There is a good input.} Suppose that there exists a good input $x_0^\star\in[k]$ such that $\{x_0^\star\}\times[a]\subseteq P_q$ for some $q$. Let $\mathcal L_{x_0^\star}$ be the set of leaves of $\pi$ that contain an input where $x_0=x_0^\star$, that is,
\[\mathcal L_{x_0^\star}=\{\ell\in\mathcal L:\exists(x_1,y_0,y_1,z,i)\text{ with }\pi(x_0^*,x_1;y_0,y_1,z,i)=\ell\}\]
and let $\overline{\mathcal L_{x_0^\star}}$ denote the complementary set of leaves, that is $\overline{\mathcal L_{x_0^\star}}=\mathcal L\setminus\mathcal L_{x_0^\star}$. We will show that $|\mathcal L_{x_0^\star}|\ge L_d^B(f)$ and that $|\overline{\mathcal L_{x_0^\star}}|\ge2k-2$. As a result, we get that
\[L(\pi)\ge|\mathcal L_{x_0^\star}|+\overline{\mathcal L_{x_0^\star}}\ge L_d^B(f)+2k-2,\]
as desired.
First, we show that $|\mathcal L_{x_0^\star}|\ge L_d^B(f)$. Let $\pi'$ be the $d$-round Bob-first subprotocol of $\pi$ given when Alice says that her input is in $P_q$ at the first round. Since $\{x_0^\star\}\times[a]\subseteq P_q$, it follows that $\pi'$ computes $f(x_1,y_1)$ on input $(x_0^\star,x_1;y_0,y_1,0,1)$ for all $x_1\in[a]$ and $y_1\in[b]$. This yields a $d$-round Bob-first protocol for computing $f$, and therefore, the number of leaves in $\pi'$ containing the input $x_0^\star$ must be at least $L_d^B(f)$. Hence, $|\mathcal L_{x_0^\star}|\ge L_d^B(f)$.
Next, we argue that $|\overline{\mathcal L_{x_0^\star}}|\ge2k-2$. Consider the set of leaves given by
\[\{\pi(x_0,x_1;y_0,y_1,z,i):i=0,x_0=y_0\in[k]\setminus\{x_0^\star\},y_1=1,z\in\{0,1\}\}.\]
If we consider the restriction of $\pi$ to the inputs $y_1=1$ and $i=0$, we can apply Lemma~20 to conclude that all $2k-2$ leaves in this set are in $\overline{\mathcal L_{x_0^\star}}$ and are pairwise distinct.
\paragraph{Case 2: No good input.} Now we consider the case where there is no good input. For each $x_0\in[k]$, we define a set $A_{x_0}\subseteq[a]$ of cardinality 2 as follows. Since $x_0\in[k]$ is not good, there exists a pair $(x_1,x'_1)\in[a]^2$ such that $\pi(x_0,x_1;y_0,y_1,z,i)\ne\pi(x_0,x'_1;y_0,y_1,z,i)$ for all $y_0,y_1,z$ and $i$. Let $A_{x_0}=\{x_1,x'_1\}$. This completes our definition of $A_{x_0}$.
We claim that the $4k$ inputs in the following set are all in pairwise distinct leaves:
\[W=\{(x_0,x_1,y_0,y_1,z,i):x_0=y_0\in[k],i=0,x_1\in A_{x_0},y_1=1,z\in\{0,1\}\}.\]
To see this, suppose that $w\ne w'$ for some $w,w'\in W$. Let $w=(x_0,x_1;y_0,y_1,z,i)$ and $w'=(x'_0,x'_1;y'_0,y'_1,z',i')$. We prove $\pi(w)\ne\pi(w')$ by considering the following three cases.
\begin{enumerate}
\item If $x_0=x'_0$ and $x_1\ne x'_1$, then we know that $\{x_1,x'_1\}=A_{x_0}$. By the construction of $A_{x_0}$, we can conclude that $\pi(w)\ne\pi(w')$.
\item If $x_0=x'_0$ and $x_1=x'_1$, then we must have $x_0=y_0=x'_0$ and $z\ne z'$ since $w\ne w'$; using Lemma~20, we conclude that $\pi(w)\ne\pi(w')$.
\item Lastly, suppose that $x_0\ne x'_0$. If we consider the restriction of $\pi$ to the inputs $y_1=1$ and $i=0$, we can apply Lemma~20 to conclude that $\pi(w)\ne\pi(w')$.
\end{enumerate}
\end{proof}
Using this lower bound, we show one can approximately compute round-$d$ complexity given an oracle that approximately computes round-$(d+1)$ complexity. We consider the following notion of approximation.
\begin{definition}
For every constant $\epsilon>0$, we say that an oracle $\mathcal O$ is a $(1+\epsilon)$-approximation of a function $\mathcal L(\cdot)$ if there exists a constant $c$ such that, for all $g$ in the domain of $\mathcal L(\cdot)$,
\[\mathcal L(g)\le\mathcal O(g)\le(1+\epsilon)\cdot\mathcal L(g)+c.\]
\end{definition}
\begin{corollary}\label{cor:approx-elimination}
Let $0<\epsilon<\frac18$. Given an oracle computing a $(1+\epsilon)$-approximation of $L_{d+1}^A(\cdot)$ and a function $f:[a]\times[b]\to\{0,1\}$, one can deterministically compute a $(1+4\epsilon)$-approximation of $L_d^B(f)$ in time $(ab)^{O(1)}$.
\end{corollary}
\begin{proof}
First, we give the reduction algorithm and then we prove its correctness. Suppose $\mathcal O$ is an oracle that computes an approximation of $L_{d+1}^A(\cdot)$ satisifying, for all functions $g$,
\[L_{d+1}^A(g)\le\mathcal O(g)\le(1+\epsilon)L_{d+1}^A(g)+O(1).\]
For a positive integer $k$, let $F_k:([k]\times[a])\times([k]\times[b]\times\{0,1\}\times\{0,1\})\to\{0,1\}$ be given by
\[F_k(x_0,x_1;y_0,y_1,z,i)=\begin{cases}\XorEq_k(x_0;y_0,z)&\text{if }i=0,\\f(x_1;y_1)&\text{if }i=1.\end{cases}\]
The reduction computes
\[v:=\max\{\mathcal O(F_k)-2(1+\epsilon)k:k\in[ab]\},\]
and outputs $v':=(v+2)/(1-2\epsilon)$. It is easy to see that this reduction runs in time $(ab)^{O(1)}$.
To prove the correctness of the reduction, we claim that
\[L_d^B(f)\le v'\le(1+4\epsilon)\cdot L_d^B(f)+O(1)\]
for all functions $f$.
From Lemma~21, we know that for all $k$
\begin{align*}
\mathcal O(F_k)-2(1+\epsilon)k
&\le(1+\epsilon)L_{d+1}^A(F_k)-2(1+\epsilon)k+O(1)\\
&\le(1+\epsilon)L_d^B(f)+O(1)
\end{align*}
so we have that
\[v'=\frac{v+2}{1-2\epsilon}\le\frac{1+\epsilon}{1-2\epsilon}\cdot L_d^B(f)+O(1)\le(1+4\epsilon)\cdot L_d^B(f)+O(1),\]
where the last inequality holds because $\epsilon<1/8$.
On the other hand, if $k=L_d^B(f)$, we have from Lemma~21 that
\begin{align*}
\mathcal O(F_k)-2(1+\epsilon)k
&\ge L_{d+1}^A(F_k)-2(1+\epsilon)k\\
&\ge\min\{4k,2k-2+L_d^B(f)\}-2(1+\epsilon)k\\
&=\min\{4L_d^B(f),3L_d^B(f)-2\}-2(1+\epsilon)L_d^B(f)\\
&=3L_d^B(f)-2-2(1+\epsilon)L_d^B(f)\\
&\ge(1-2\epsilon)L_d^B(F_k)-2.
\end{align*}
Since $k=L_d^B(f)\le ab$, we conclude that $v'=(v+2)/(1-2\epsilon)\ge L_d^B(F_k)$.
\end{proof}
Combining Corollary~23 with the hardness result for the $d=3$ case in Theorem~18, we get that computing $L_d$ is NP-hard (under a polynomial-time truth-table reduction).

\begin{corollary}
For any integer $d\ge3$, there exists an $\epsilon>0$ such that given access to an oracle that computes a $(1+\epsilon)$-approximation of $L_d^A$, one can compute any language in $\mathrm{NP}$ in polynomial time.
\end{corollary}

\begin{thebibliography}{65}
\bibitem{ref1} Scott Aaronson and Avi Wigderson. Algebrization: A new barrier in complexity theory. ACM Transactions on Computation Theory, 1(1):2, 2009.
\bibitem{ref2} Miklós Ajtai. $\Sigma^1_1$-formulae on finite structures. Annals of pure and applied logic, 24(1):1-48, 1983.
\bibitem{ref3} Miklós Ajtai and Michael Ben-Or. A Theorem on Probabilistic Constant Depth Computations. In Proceedings of the Symposium on Theory of Computing (STOC), pages 471-474, 1984.
\bibitem{ref4} Eric Allender. The new complexity landscape around circuit minimization. In Alberto Leporati, Carlos Martín-Vide, Dana Shapira, and Claudio Zandron, editors, Language and Automata Theory and Applications (LATA), volume 12038 of Lecture Notes in Computer Science, pages 3-16. Springer, 2020.
\bibitem{ref5} Eric Allender and Bireswar Das. Zero knowledge and circuit minimization. In International Symposium on Mathematical Foundations of Computer Science (MFCS), pages 25-32, 2014.
\bibitem{ref6} Eric Allender, Joshua A. Grochow, Dieter van Melkebeek, Cristopher Moore, and Andrew Morgan. Minimum circuit size, graph isomorphism, and related problems. SIAM Journal on Computing, 47(4):1339-1372, 2018.
\bibitem{ref7} Eric Allender, Lisa Hellerstein, Paul McCabe, Toniann Pitassi, and Michael E. Saks. Minimizing disjunctive normal form formulas and AC$^0$ circuits given a truth table. SIAM Journal on Computing, 38(1):63-84, 2008.
\bibitem{ref8} Eric Allender and Shuichi Hirahara. New insights on the (non-)hardness of circuit minimization and related problems. In International Symposium on Mathematical Foundations of Computer Science (MFCS), pages 54:1-54:14, 2017.
\bibitem{ref9} Eric Allender, Dhiraj Holden, and Valentine Kabanets. The minimum oracle circuit size problem. Computational Complexity, 26(2):469-496, 2017.
\bibitem{ref10} Eric Allender, Rahul Ilango, and Neekon Vafa. The non-hardness of approximating circuit size. In International Computer Science Symposium in Russia (CSR), pages 13-24, 2019.
\bibitem{ref11} Eric Allender and Michal Koucky. Amplifying lower bounds by means of self-reducibility. Journal of the ACM, 57(3):1-36, 2010.
\bibitem{ref12} Eric Allender, Michal Koucky, Detlef Ronneburger, and Sambuddha Roy. The pervasive reach of resource-bounded Kolmogorov complexity in computational complexity theory. Journal of Computer and System Sciences, 77(1):14-40, 2011.
\bibitem{ref13} Noga Alon and Joel H Spencer. The probabilistic method. John Wiley \& Sons, 2016.
\bibitem{ref14} Theodore Baker, John Gill, and Robert Solovay. Relativizations of the p=?np question. SIAM Journal on computing, 4(4):431-442, 1975.
\bibitem{ref15} Boaz Barak, Mark Braverman, Xi Chen, and Anup Rao. How to compress interactive communication. SIAM Journal on Computing, 42(3):1327-1363, 2013.
\bibitem{ref16} Avrim L. Blum and Ronald L. Rivest. Training a 3-node neural network is np-complete. Neural Networks, 5(1):117-127, 1992.
\bibitem{ref17} Joan Boyar, Philip Matthews, and René Peralta. On the shortest linear straight-line program for computing linear forms. In International Symposium on Mathematical Foundations of Computer Science (MFCS), pages 168-179, 2008.
\bibitem{ref18} Marco L. Carmosino, Russell Impagliazzo, Valentine Kabanets, and Antonina Kolokolova. Learning algorithms from natural proofs. In Conference on Computational Complexity (CCC), pages 10:1-10:24, 2016.
\bibitem{ref19} Sebastian Lukas Arne Czort. The complexity of minimizing disjunctive normal form formulas. Master’s thesis, University of Aarhus, 1999.
\bibitem{ref20} Uriel Feige and Joe Kilian. Zero knowledge and the chromatic number. Journal of Computer and System Sciences, 57(2):187-199, 1998.
\bibitem{ref21} Vitaly Feldman. Hardness of approximate two-level logic minimization and PAC learning with membership queries. In Symposium on Theory of Computing (STOC), pages 363-372, 2006.
\bibitem{ref22} Alexander Golovnev, Rahul Ilango, Russell Impagliazzo, Valentine Kabanets, Antonina Kolokolova, and Avishay Tal. ac$^0$[p] lower bounds against mcsp via the coin problem. In Colloquium on Automata, Languages, and Programming (ICALP), volume 132, page 66, 2019.
\bibitem{ref23} Thomas R. Hancock, Tao Jiang, Ming Li, and John Tromp. Lower bounds on learning decision lists and trees. Information and Computation, 126(2):114-122, 1996.
\bibitem{ref24} Johan Hastad. Clique is hard to approximate within n\textasciicircum{}1-$\varepsilon$. In Symposium on Foundations of Computer Science (FOCS), pages 627-636, 1996.
\bibitem{ref25} Shuichi Hirahara, Igor Carboni Oliveira, and Rahul Santhanam. NP-hardness of minimum circuit size problem for OR-AND-MOD circuits. In Computational Complexity Conference (CCC), pages 5:1-5:31, 2018.
\bibitem{ref26} Shuichi Hirahara and Osamu Watanabe. Limits of minimum circuit size problem as oracle. In Conference on Computational Complexity (CCC), pages 18:1-18:20, 2016.
\bibitem{ref27} John M. Hitchcock and Aduri Pavan. On the NP-completeness of the minimum circuit size problem. In Foundation of Software Technology and Theoretical Computer Science (FSTTCS), pages 236-245, 2015.
\bibitem{ref28} Rahul Ilango. Approaching MCSP from above and below: Hardness for a conditional variant and ac$^0$[p]. In Thomas Vidick, editor, Innovations in Theoretical Computer Science Conference (ITCS), volume 151 of LIPIcs, pages 34:1-34:26. Schloss Dagstuhl - Leibniz-Zentrum für Informatik, 2020.
\bibitem{ref29} Rahul Ilango. Constant depth formula and partial function versions of MCSP are hard. In Proceedings of the Symposium on Foundations of Computer Science (FOCS), pages 424-433. IEEE, 2020.
\bibitem{ref30} Rahul Ilango, Bruno Loff, and Igor Carboni Oliveira. Np-hardness of circuit minimization for multi-output functions. In Shubhangi Saraf, editor, Computational Complexity Conference (CCC), volume 169 of LIPIcs, pages 22:1-22:36. Schloss Dagstuhl - Leibniz-Zentrum für Informatik, 2020.
\bibitem{ref31} Russell Impagliazzo, Valentine Kabanets, and Ilya Volkovich. The Power of Natural Properties as Oracles. In Computational Complexity Conference (CCC), volume 102, pages 7:1-7:20, 2018.
\bibitem{ref32} J. Stephen Judd. Learning in networks is hard. In International Conference on Neural Networks (ICNN), volume 2, pages 685-692, 1987.
\bibitem{ref33} Valentine Kabanets and Jin-yi Cai. Circuit minimization problem. In Symposium on Theory of Computing (STOC), pages 73-79, 2000.
\bibitem{ref34} Mauricio Karchmer, Eyal Kushilevitz, and Noam Nisan. Fractional covers and communication complexity. SIAM Journal on Discrete Mathematics, 8(1):76-92, 1995.
\bibitem{ref35} Mauricio Karchmer and Avi Wigderson. Monotone circuits for connectivity require super-logarithmic depth. In Proceedings of the Twentieth Annual ACM Symposium on Theory of Computing, STOC '88, page 539–550. Association for Computing Machinery, 1988.
\bibitem{ref36} Subhash Khot and Rishi Saket. Hardness of minimizing and learning DNF expressions. In Symposium on Foundations of Computer Science (FOCS), pages 231-240, 2008.
\bibitem{ref37} Jan Krajícek. Forcing with Random Variables and Proof Complexity. Cambridge University Press, 2011.
\bibitem{ref38} Eyal Kushilevitz and Noam Nisan. Communication Complexity. Cambridge University Press, 1997.
\bibitem{ref39} Eyal Kushilevitz and Enav Weinreb. On the complexity of communication complexity. In Symposium on Theory of Computing (STOC), pages 465-474, 2009.
\bibitem{ref40} Yanyi Liu and Rafael Pass. On one-way functions and kolmogorov complexity. arXiv preprint, 2020. URL: http://arxiv.org/abs/2009.11514.
\bibitem{ref41} L. Lovasz and M. Saks. Lattices, mobius functions and communications complexity. In Symposium on Foundations of Computer Science (FOCS), SFCS '88, page 81–90, USA, 1988. IEEE Computer Society.
\bibitem{ref42} Carsten Lund and Mihalis Yannakakis. On the hardness of approximating minimization problems. Journal of the ACM, 41(5):960-981, 1994.
\bibitem{ref43} William J. Masek. Some NP-complete set covering problems. Unpublished Manuscript, 1979.
\bibitem{ref44} Dylan M. McKay, Cody D. Murray, and R. Ryan Williams. Weak lower bounds on resource-bounded compression imply strong separations of complexity classes. In Symposium on Theory of Computing (STOC), STOC 2019, page 1215–1225, New York, NY, USA, 2019. Association for Computing Machinery.
\bibitem{ref45} Dylan M. McKay, Cody D. Murray, and R. Ryan Williams. Weak lower bounds on resource-bounded compression imply strong separations of complexity classes. In Symposium on Theory of Computing (STOC), 2019.
\bibitem{ref46} Peter Bro Miltersen, Noam Nisan, Shmuel Safra, and Avi Wigderson. On data structures and asymmetric communication complexity. Journal of Computer and System Sciences, 57(1):37-49, 1998.
\bibitem{ref47} Moritz Müller and Ján Pich. Feasibly constructive proofs of succinct weak circuit lower bounds. Annals of Pure and Applied Logic, 171(2), 2020.
\bibitem{ref48} Cody D. Murray and Richard Ryan Williams. On the (non) NP-hardness of computing circuit complexity. In Conference on Computational Complexity (CCC), pages 365-380, 2015.
\bibitem{ref49} Joseph Naor and Moni Naor. Small-bias probability spaces: Efficient constructions and applications. SIAM journal on computing, 22(4):838-856, 1993.
\bibitem{ref50} Moni Naor and Omer Reingold. Number-theoretic constructions of efficient pseudo-random functions. Journal of the ACM, 51(2):231-262, 2004.
\bibitem{ref51} Moni Naor, Leonard J. Schulman, and Aravind Srinivasan. Splitters and Near-Optimal Derandomization. In Symposium on Foundations of Computer Science (FOCS), pages 182-191, 1995.
\bibitem{ref52} Noam Nisan. Pseudorandom bits for constant depth circuits. Combinatorica, 11(1):63-70, 1991.
\bibitem{ref53} Igor Carboni Oliveira, Ján Pich, and Rahul Santhanam. Hardness magnification near state-of-the-art lower bounds. In Conference on Computational Complexity (CCC), 2019.
\bibitem{ref54} Igor Carboni Oliveira and Rahul Santhanam. Conspiracies between learning algorithms, circuit lower bounds, and pseudorandomness. In Computational Complexity Conference (CCC), pages 18:1-18:49, 2017.
\bibitem{ref55} Igor Carboni Oliveira and Rahul Santhanam. Hardness magnification for natural problems. In Symposium on Foundations of Computer Science (FOCS), pages 65-76, 2018.
\bibitem{ref56} Denis Pankratov. Direct sum questions in classical communication complexity. Master’s thesis, University of Chicago, 2012.
\bibitem{ref57} Leonard Pitt and Leslie G. Valiant. Computational limitations on learning from examples. Journal of the ACM, 35(4):965-984, 1988.
\bibitem{ref58} Anup Rao and Amir Yehudayoff. Communication Complexity: and Applications. Cambridge University Press, 2020.
\bibitem{ref59} Alexander A Razborov. On submodular complexity measures. Boolean Function Complexity,(M. Paterson, Ed.), pages 76-83, 1992.
\bibitem{ref60} Alexander A Razborov and Steven Rudich. Natural proofs. Journal of Computer and System Sciences, 55(1):24-35, 1997.
\bibitem{ref61} Alexander A. Razborov and Steven Rudich. Natural proofs. Journal of Computer and System Sciences, 55(1):24-35, 1997.
\bibitem{ref62} Pranab Sen and Srinivasan Venkatesh. Lower bounds for predecessor searching in the cell probe model. Journal of Computer and System Sciences, 74(3):364-385, 2008.
\bibitem{ref63} Boris A Trakhtenbrot. A survey of Russian approaches to perebor (brute-force search) algorithms. Annals of the History of Computing, 6(4):384-400, 1984.
\bibitem{ref64} Christopher Umans, Tiziano Villa, and Alberto L. Sangiovanni-Vincentelli. Complexity of two-level logic minimization. IEEE Transactions on CAD of Integrated Circuits and Systems, 25(7):1230-1246, 2006.
\bibitem{ref65} David Zuckerman. Linear Degree Extractors and the Inapproximability of Max Clique and Chromatic Number. Theory of Computing, 3(1):103-128, 2007.
\end{thebibliography}
\end{document}
